% Généré par tools/build.py — ne pas éditer à la main
\section{Proportions et évolutions}

\fichehead{mar.~2~févr.~2027}{Proportions et pourcentages\,; proportion de proportion}
\begin{fiche}
\qitem{Échauffement · Flash}{PRO}{Calculer $15\,\%$ de $240$.}{$0{,}15\times240=36$.}{}
\qitem{Cœur · Méthode}{PRO}{$30\,\%$ des élèves sont internes et $40\,\%$ des internes pratiquent la danse. Quelle proportion des élèves sont des internes qui dansent\,?}{$0{,}3\times0{,}4=0{,}12=12\,\%$~: proportion de proportion (on multiplie).}{}
\qitem{Cœur + · Problème}{PRO}{Un lycée compte $800$ élèves dont $55\,\%$ de filles\,; $20\,\%$ des filles font du latin. Combien de filles latinistes\,?}{$800\times0{,}55=440$ filles\,; $440\times0{,}20=88$.}{}
\qitem{Retour · Flash}{QP}{$3$ cahiers coûtent $4{,}50$ €. Prix de $7$ cahiers\,?}{$1{,}50\times7=10{,}50$ €.}{}
\qitem{Fil rouge · Vrai/faux}{EQ}{Vrai ou faux\,? Justifier. Le nombre $-2$ est solution de l'équation $3x+1=-5$.}{Vrai. $3\times(-2)+1=-6+1=-5$.}{}
\end{fiche}

\fichehead{jeu.~4~févr.~2027}{Variation absolue, coefficient multiplicateur, taux d'évolution}
\begin{fiche}
\qitem{Échauffement · Flash}{PRO}{Un prix passe de $80$ Ar à $92$ Ar. Quelle est la variation absolue\,?}{$92-80=12$ Ar.}{}
\qitem{Cœur · Méthode}{PRO}{Pour la même évolution, donner le taux d'évolution et le coefficient multiplicateur.}{$t=\dfrac{12}{80}=0{,}15=15\,\%$\,; $CM=1+t=1{,}15$.}{}
\qitem{Cœur + · QCM}{PRO}{Une baisse de $20\,\%$ correspond au coefficient multiplicateur~:\choix{$0{,}2$}{$-0{,}8$}{$1{,}2$}{$0{,}8$}}{D.\ $0{,}8$}{$CM=1+t=1-0{,}20$.}
\qitem{Retour · Flash}{STA}{Températures $-2,\,0,\,3,\,5,\,9$ (en °C). Moyenne et étendue\,?}{Moyenne $3$ °C\,; étendue $9-(-2)=11$ °C.}{}
\qitem{Fil rouge · Flash}{SOL}{Calculer le volume et l'aire totale d'un cube d'arête $3$ cm.}{$V=3^3=27$ cm$^3$\,; $A=6\times3^2=54$ cm$^2$.}{}
\end{fiche}

\fichehead{lun.~8~févr.~2027}{Évolutions successives}
\begin{fiche}
\qitem{Échauffement · Flash}{CAL}{Calculer $1{,}2\times0{,}9$.}{$1{,}08$.}{}
\qitem{Cœur · Méthode}{PRO}{Une hausse de $20\,\%$ est suivie d'une baisse de $10\,\%$. Quel est le taux global\,?}{$1{,}2\times0{,}9=1{,}08$~: hausse globale de $8\,\%$.}{Les taux ne s'additionnent pas\,; ce sont les coefficients qui se multiplient.}
\qitem{Cœur + · Vrai/faux}{PRO}{Vrai ou faux\,? Justifier. Une hausse de $50\,\%$ suivie d'une baisse de $50\,\%$ ramène à la valeur initiale.}{Faux. $1{,}5\times0{,}5=0{,}75$~: baisse globale de $25\,\%$.}{}
\qitem{Retour · Flash}{SUB}{Calculer $x^2-x$ pour $x=\dfrac12$.}{$-\dfrac{1}{4}$.}{}
\qitem{Fil rouge · Rédiger}{DEM}{$ABCD$ est un parallélogramme tel que $AC=BD$. Quelle est sa nature\,? Justifier en une phrase.}{Rectangle~: un parallélogramme dont les diagonales ont même longueur est un rectangle.}{}
\end{fiche}

\fichehead{mar.~9~févr.~2027}{Évolution réciproque\,; proportion et évolution}
\begin{fiche}
\qitem{Échauffement · Flash}{CAL}{Calculer $\dfrac{1}{1{,}25}$.}{$0{,}8$.}{}
\qitem{Cœur · Méthode}{PRO}{Après une hausse de $25\,\%$, quelle baisse ramène à la valeur initiale\,?}{$CM_{\text{réc}}=\dfrac{1}{1{,}25}=0{,}8$~: baisse de $20\,\%$ (et non $25\,\%$).}{Piège~: croire qu'une hausse de $25\,\%$ se compense par une baisse de $25\,\%$.}
\qitem{Cœur + · Interpréter}{PRO}{Sur $250$ candidats, $150$ sont admis\,; l'an dernier il y avait $200$ candidats. Calculer la proportion d'admis puis le taux d'évolution du nombre de candidats.}{Proportion $\dfrac{150}{250}=60\,\%$ (partie sur tout)\,; évolution $\dfrac{250-200}{200}=25\,\%$ (variation relative).}{}
\qitem{Retour · Tableau}{IMG}{Tableau de valeurs de $f$~:\par\smallskip\renewcommand{\arraystretch}{1.15}\begin{tabular}{|c|c|c|c|c|}\hline
$x$ & 0 & 1 & 2 & 3 \\ \hline
$f(x)$ & 4 & 1 & 0 & 1 \\ \hline
\end{tabular}\par Image de $2$\,? Antécédents de $1$\,?}{$f(2)=0$\,; antécédents de $1$~: $1$ et $3$.}{}
\qitem{Fil rouge · QCM}{REP}{Sur le globe terrestre, deux points d'un même parallèle ont~:\choix{la même latitude}{la même longitude}{la même distance à Greenwich}{la même altitude}}{A.\ la même latitude}{Latitude~: angle nord/sud par rapport à l'équateur\,; longitude~: angle est/ouest par rapport au méridien de Greenwich.}
\end{fiche}
